\documentclass[11pt]{article}
\usepackage[margin=1in]{geometry}
\usepackage{amsmath, amssymb}
\usepackage{enumitem}

\begin{document}

\begin{center}
    {\Large \textbf{Summary: Absorbing Markov Chain Analysis}}\\[0.5em]
\end{center}

This problem illustrates the standard workflow for analyzing a finite Markov chain with transient states and multiple closed classes.  
The key steps and concepts are summarized below so the method can be applied to similar problems.

\section*{(a) Canonical Form of the Transition Matrix}

The first step is to classify states as:
\begin{itemize}[leftmargin=2em]
    \item \textbf{Transient states}: states that may be left permanently.
    \item \textbf{Closed (absorbing) classes}: once entered, the chain never leaves.
\end{itemize}

After reordering states so that all transient states come first, the transition matrix takes the \textbf{canonical form}:
\[
P =
\begin{pmatrix}
Q & R \\
0 & S
\end{pmatrix},
\]
where:
\begin{itemize}[leftmargin=2em]
    \item $Q$ contains transitions among transient states,
    \item $R$ contains transitions from transient states to closed states,
    \item $S$ contains transitions within closed classes.
\end{itemize}

This reordering does not change the chain, but exposes its long-term structure.

\section*{(b) Fundamental Matrix}

The \textbf{fundamental matrix} is defined as
\[
F = (I - Q)^{-1}.
\]

Its interpretation is central:
\begin{quote}
$F_{ij}$ is the expected number of visits to transient state $j$ given the chain starts in transient state $i$.
\end{quote}

Equivalently,
\[
F = I + Q + Q^2 + Q^3 + \cdots,
\]
which sums expected visits over all time before absorption.

\section*{(c) Expected Time to Absorption}

The expected number of steps before entering a closed class is obtained by summing the entries of the corresponding row of $F$:
\[
\mathbb{E}_i[T_{\text{abs}}] = \sum_j F_{ij}.
\]

Each visit to a transient state corresponds to one unit of time, so row sums represent total expected time spent in transient states.

\section*{(d) Absorption Probabilities}

Define
\[
B = FR.
\]

Then:
\begin{quote}
$B_{ij}$ is the probability that the chain is eventually absorbed in closed state (or class) $j$, given it started in transient state $i$.
\end{quote}

This formula follows from linearity of expectation:
\begin{itemize}[leftmargin=2em]
    \item $F_{ik}$ counts how many times state $k$ is visited on average,
    \item $R_{kj}$ gives the probability of exiting from $k$ to $j$,
    \item Multiplying and summing accounts for all possible exit times.
\end{itemize}

\section*{(e) Limiting Distribution}

If the chain has multiple closed classes, there is \textbf{no unique global limiting distribution}.  
Instead, the limit depends on the starting state.

The limiting distribution is constructed as follows:
\begin{enumerate}[leftmargin=2em]
    \item Transient states receive probability $0$ in the limit.
    \item Absorption probabilities (from $B$) determine which closed class is entered.
    \item Within each closed class, the chain converges to that class’s \textbf{stationary distribution}.
\end{enumerate}

The final limiting distribution is a weighted mixture of the stationary distributions of the closed classes, with weights given by absorption probabilities.

\section*{Key Takeaway}

\begin{quote}
\textit{Absorbing Markov chain analysis separates behavior into three phases:  
(1) movement among transient states,  
(2) probabilities of entering closed classes, and  
(3) long-run behavior within those classes.}
\end{quote}

This framework applies broadly to Markov chains with transient behavior followed by absorption.

\newpage
\section*{Estimating and Smoothing Markov Transition Probabilities}

Consider a time-homogeneous Markov chain with finite state space
\[
\mathcal{S} = \{1,2,\dots,K\},
\]
and let \(n_{ij}\) denote the observed number of transitions from state \(i\) to state \(j\).
Define the row totals
\[
n_{i\cdot} = \sum_{j=1}^K n_{ij}.
\]

\subsection*{Maximum Likelihood Estimation (MLE)}

Conditioning on being in state \(i\), the next-state transitions follow a multinomial distribution.
The maximum likelihood estimator of the transition probabilities is obtained by row normalization:
\[
\hat p_{ij}
=
\Pr(X_{n+1}=j \mid X_n=i)
=
\frac{n_{ij}}{n_{i\cdot}},
\quad j=1,\dots,K.
\]

This estimator is simple and unbiased asymptotically, but can be unstable when \(n_{i\cdot}\) is small.

\subsection*{Limitations of the MLE}

If a particular state \(i\) is rarely observed (small \(n_{i\cdot}\)), then:
\begin{itemize}
  \item Estimates have high variance.
  \item Zero counts produce zero probability estimates, even if the true transition probability is nonzero.
\end{itemize}
These issues motivate smoothing or regularization.

\subsection*{Laplace (Add-One) Smoothing}

Laplace smoothing adjusts the counts by adding a constant (typically 1) to each possible transition:
\[
\tilde n_{ij} = n_{ij} + 1.
\]

The smoothed transition probabilities are then:
\[
\tilde p_{ij}
=
\frac{n_{ij} + 1}{n_{i\cdot} + K}.
\]

Laplace smoothing guarantees strictly positive probabilities but does not incorporate external structural information.

\subsection*{Dirichlet Smoothing with an Empirical Bayes Prior}

A more principled approach places a Dirichlet prior on each transition row:
\[
(p_{i1},\dots,p_{iK}) \sim \text{Dirichlet}(\alpha_1,\dots,\alpha_K).
\]

Let the prior mean be \(m = (m_1,\dots,m_K)\) with \(\sum_j m_j = 1\), and let \(\alpha > 0\) denote the prior strength.
Then:
\[
\alpha_j = \alpha m_j.
\]

After observing transition counts \(n_{ij}\), the posterior distribution is:
\[
(p_{i1},\dots,p_{iK}) \mid \text{data}
\sim
\text{Dirichlet}(\alpha_1 + n_{i1},\dots,\alpha_K + n_{iK}).
\]

The posterior mean (used as a smoothed estimator) is:
\[
\tilde p_{ij}
=
\frac{\alpha m_j + n_{ij}}{\alpha + n_{i\cdot}}.
\]

\subsection*{Interpretation}

\begin{itemize}
  \item \(\alpha\) controls the influence of the prior relative to the data.
  \item The estimator is a convex combination of the prior mean and the MLE.
  \item Empirical Bayes chooses \(m\) from pooled or related data, allowing information sharing across states.
\end{itemize}

\end{document}
